\subsection{Định tuyến theo chất lượng liên kết}
FANET có tính di động cao và topology thay đổi nhanh~\cite{bekmezci}, nên các
khảo sát về mạng UAV~\cite{gupta,hayat} và về định tuyến
FANET~\cite{arafat,lakew} đều xem độ bền liên kết là yếu tố quyết định.
OLSR~\cite{rfc3626} chọn tuyến theo số hop: một liên kết được coi là dùng được
chỉ cần HELLO đi qua hai chiều, mức tổn thất gói không vào chi phí tuyến.
OLSRv2~\cite{rfc7181} thêm khung metric nhưng bản OLSR trong ns-3 theo RFC 3626.
ETX~\cite{decouto} và ETT~\cite{draves} thay số hop bằng chi phí phản ánh tổn
thất và băng thông; olsrd~\cite{olsrd} đưa ETX vào OLSR và là baseline phân tán
tự nhiên cho mọi phương pháp định tuyến theo chất lượng liên kết.

\subsection{Dự báo chất lượng liên kết và GNN}
Ước lượng chất lượng liên kết là một hướng riêng, từ các bộ ước lượng thống kê
gói~\cite{baccour} đến các mô hình học máy~\cite{cerar}. GNN học biểu diễn nút
bằng truyền thông điệp lặp lại: GCN~\cite{kipf}, GraphSAGE~\cite{graphsage} với
tổng hợp lân cận áp dụng được cho đồ thị chưa thấy, GAT~\cite{gat} gán trọng số
chú ý. Mạng UAV tự nhiên là đồ thị biến thiên theo thời gian nên GNN là ứng viên
hợp lý. Trong mạng có dây, RouteNet~\cite{routenet} dùng GNN dự đoán trễ và mất
gói từ topology và tuyến. Trong mạng UAV, Bai và cộng sự~\cite{bai2025} dự đoán
liên kết bằng GNN động và lấy chính GraphSAGE cùng node2vec làm đối chứng;
GNNPPOR~\cite{gnnppor} ghép GNN với PPO để chọn tuyến trong FANET; một khảo sát
gần đây~\cite{drones2026} tổng hợp các hướng dùng GNN cho hệ thống UAV. Trong
FANET, QMR~\cite{qmr} dùng Q-learning chọn bước kế tiếp theo nhiều mục tiêu, và
XGBoost~\cite{xgboost} là baseline bảng mạnh cần có khi đánh giá GNN.
Predictive-OLSR~\cite{rosati2013} dùng GPS và vận tốc tương đối để điều chỉnh
trọng số ETX; OLSR-LCN~\cite{liu2025} kết hợp tuổi thọ liên kết, nhiễu và tải
nút. Nghĩa là \emph{dùng dự báo, vị trí hoặc GNN cho định tuyến UAV không còn
là ý tưởng mới}; câu hỏi còn mở là lợi ích thực tế của nó khi đo bằng dịch vụ.

\subsection{Khoảng cách giữa độ chính xác và lợi ích dịch vụ}
Khảo sát của Cerar và cộng sự~\cite{cerar} nhận xét rằng phần lớn mô hình ước
lượng chất lượng liên kết được đánh giá ngoại tuyến bằng độ chính xác, ít khi
được đặt vào vòng định tuyến để đo lợi ích dịch vụ. Shi và cộng
sự~\cite{shi2024} cho thấy nhiều họ mô hình học máy cùng chạm một trần độ chính
xác do tính ngẫu nhiên của kênh. Ở các lĩnh vực khác, khoảng cách này đã được
ghi nhận rõ: trong lái tự động, thứ hạng của các bộ dự báo chuyển động thay đổi
đáng kể khi chuyển từ đánh giá hở sang đánh giá vòng kín~\cite{closedloop2025};
trong truy hồi thông tin, phần lớn lợi ích báo cáo của mô hình nơ-ron biến mất
khi baseline được chỉnh đúng~\cite{lin2018}; trong học máy cho phương trình đạo
hàm riêng, 79\% bài so sánh với một baseline yếu~\cite{mcgreivy2024}. Đăng ký
trước~\cite{nosek} là biện pháp trực tiếp chống việc chọn lại cấu hình sau khi
thấy kết quả.

\subsection{Điều khiển vị trí relay}
Hướng tối ưu quỹ đạo relay~\cite{zeng2016} điều khiển vị trí UAV để tối đa
thông lượng, thường với mô hình kênh giải tích và giả định tuyến đi qua chuỗi
relay mà bộ điều khiển chọn. Mạng UAV định nghĩa bằng phần mềm~\cite{abir} dùng
bộ điều khiển có tầm nhìn toàn mạng, với nhược điểm thường được nêu là chi phí
báo hiệu, độ trễ điều khiển và khả năng mở rộng. Nghiên cứu này không giải bài
toán quỹ đạo tối ưu; nó giữ một luật bố trí cách đều cố định để tách riêng
\emph{thời điểm kích hoạt}, và đo thêm một giới hạn hình học để biết luật đó
còn cách mức tốt nhất khả thi bao xa.

\subsection{Khoảng trống được nhắm tới}
Chưa có công trình nào trong mạng UAV, theo hiểu biết của nhóm, (i) cho cùng
một bộ dự báo đã khóa điều khiển hai hành động khác nhau trong cùng simulator,
(ii) đăng ký trước endpoint và thứ tự kiểm định rồi báo cáo cả kết quả bất lợi,
và (iii) đối chiếu thứ hạng theo sai số dự báo với thứ hạng theo dịch vụ trên
cùng tập trace. Đó là ba điểm bài báo này nhắm tới.
